\section{Application, Deployment and Tracking}\label{sec:app}
\textbf{Interface design.} The interface was designed first in Google Stitch~\cite{stitch}. A first prompt described four workspaces with a sidebar, input and corruption cards, result panels and a system-information area; Stitch returned four screens and a design system (indigo, Inter, white rounded cards). Using them, we found that the controls and Run button sat below the fold, so a run needed scrolling down to choose options and back up to see the result. A second prompt asked for a layout where everything for one run fits on one laptop screen (narrow control panel with a pinned Run button, large results on the right). Both rounds and the project page are shown in Fig.~\ref{fig:stitch}; the prompts are in \code{stitch\_design/PROMPTS.md}. The exported HTML holds invented data (GPU model, LPIPS), so only its structure, colours and typography were re-implemented as React components that show real data.

\begin{figure*}[t]
\centering
\includegraphics[width=0.49\textwidth]{stitch_montage.pdf}\hfill
\includegraphics[width=0.49\textwidth]{stitch_project.png}
\caption{Left: Stitch designs of the first prompt (top) and the layout revision (bottom). Right: the Stitch project page with both rounds and the prompt list.}\label{fig:stitch}
\end{figure*}

\textbf{Front end and back end.} The React~18 / Vite / Tailwind front end has four workspaces (Fig.~\ref{fig:screens}): the user uploads an image, picks a built-in test sample or (for sketches) uses the webcam, chooses a corruption and severity (or custom sliders and a seed) or none, and presses Run. Results show the model input, output, absolute error map, inference time and corruption settings; hard routing adds the four classifier probabilities, the predicted corruption, the selected expert and oracle/predicted mode; the mixture adds the four weights, a stacked contribution bar and the dominant branch; sketches show photograph and sketch side by side with a download button. A browser measurement showed that at $1366\times768$ and $1920\times1080$ the Run button is always visible and the page does not scroll. The FastAPI service loads the seven ONNX models once and offers \code{/api/health}, \code{/api/restore/universal}, \code{/hard}, \code{/soft} and \code{/api/sketch}. It validates uploads (type, 10~MB and 40~megapixel limits, decodability, style and slider ranges), preprocesses exactly as in training, applies corruptions with a NumPy port of the training code, runs ONNX Runtime on the CPU and returns images, routing information and timings; a missing model gives a clear 503 message instead of a crash.

\begin{figure*}[t]
\centering
\includegraphics[width=0.36\textwidth]{app_universal.png}\hfill
\includegraphics[width=0.36\textwidth]{app_hard.png}\\[2pt]
\includegraphics[width=0.36\textwidth]{app_soft.png}\hfill
\includegraphics[width=0.36\textwidth]{app_sketch.png}
\caption{The deployed application ($1366\times768$): universal, hard-routed, soft mixture and face-to-sketch workspaces after a run on a built-in sample.}\label{fig:screens}
\end{figure*}

\textbf{ONNX.} All inference models were exported (opset 17) and compared with PyTorch on a random batch; the maximum absolute differences are all below $2.1\times10^{-6}$ (Table~\ref{tab:onnx}). The soft mixture (gate, experts, identity and mixing) is one ONNX graph returning the image and the four weights; only the cGAN generator is exported. The reported test results were also re-measured using only the deployed ONNX files, the back end's corruptions and a different SSIM implementation; SSIM agreed to four decimals on the first 300 test entries.

\begin{table}[t]
\centering
\caption{ONNX files: size, maximum absolute difference to PyTorch, CPU latency per image.}\label{tab:onnx}
\small
\adjustbox{max width=\columnwidth}{\begin{tabular}{lrrr}
\toprule
Model & Size (MB) & Max $|$PyTorch$-$ONNX$|$ & Latency (ms) \\
\midrule
Task 1 universal autoencoder & 16.1 & 5.4e-7 & 12.4 \\
Task 2 classifier & 1.2 & 2.0e-6 & 0.3 \\
Task 2 specialist (salt-and-pepper) & 16.1 & 5.7e-7 & 15.3 \\
Task 2 specialist (blur) & 16.1 & 4.5e-7 & 12.4 \\
Task 2 specialist (occlusion) & 16.1 & 5.4e-7 & 15.3 \\
Task 3 soft mixture (whole pipeline) & 49.4 & 2.4e-7 & 39.8 \\
Task 4 generator & 168.5 & 1.4e-6 & 21.1 \\
\bottomrule
\end{tabular}
}
\end{table}

\textbf{Deployment.} Two containers are built from the repository, the back end (Python, FastAPI, ONNX Runtime) and the front end (nginx serving the build and proxying \code{/api}). The model files are mounted from \code{./onnx\_models}, are not stored in Git and are fetched by \code{scripts/download\_models.py} from the \code{models-v1} release of the repository. After cloning, \code{docker compose up --build} starts the application at \code{http://localhost:3000}; an optional profile starts MLflow on port 5000. A smoke test exercises every endpoint, option and error path through the nginx address (54 checks), and the back end has 38 unit tests. \code{README.md} and \code{KAGGLE\_GUIDE.md} give the complete steps to run, retrain and evaluate.

\textbf{Experiment tracking.} Runs, Optuna trials, hyper-parameters, per-epoch losses and metrics, checkpoints and visual outputs (Task~4 sample grids every ten epochs) are recorded with MLflow in five experiments; the MLflow database and Optuna studies are in the repository (Fig.~\ref{fig:mlflow}).

\begin{figure*}[t]
\centering
\includegraphics[width=0.36\textwidth]{mlflow_experiments.png}\hfill
\includegraphics[width=0.36\textwidth]{mlflow_run.png}
\caption{MLflow: the experiments with their runs (left) and the metrics of one run (right).}\label{fig:mlflow}
\end{figure*}
